\title{The Era of 1-bit LLMs: All Large Language Models are in 1.58 Bits}

\begin{abstract}
Emerging research exemplified by BitNet~\cite{bitnet} signals the advent of 1-bit Large Language Models (LLMs). This study presents \textbf{\text{BitNet b1.58}}, a 1-bit LLM version where every model parameter (weight) is restricted to the ternary set \{-1, 0, 1\}. This approach achieves parity with its full-precision Transformer LLM counterparts (using FP16 or BF16) of identical model size and trained with the same quantity of tokens, both in perplexity and in downstream task performance. Moreover, it yields marked improvements in efficiency, exhibiting substantial reductions in latency, memory footprint, throughput, and power demands. Crucially, the proposed 1.58-bit LLM articulates a novel scaling law alongside methodological guidelines for constructing future LLMs that balance high performance with economical resource utilization. Additionally, it supports a reimagined computational framework and lays foundational principles for future hardware tailored specifically to 1-bit LLM architectures.
\end{abstract}

\section{The Era of 1-bit LLMs}

Recent advancements in AI have led to exponential growth in both the scale and ability of Large Language Models (LLMs). While these models achieve exceptional results across an array of natural language processing tasks, their increasing magnitude introduces barriers to practical deployment and exacerbates environmental as well as financial costs primarily because of heightened energy usage.

A strategy to mitigate these limitations is post-training quantization, which generates low-bitwidth models tailored for inference~\cite{smoothquant, gptq, quip, quip_sharp}. By reducing the numerical precision of both weights and activations, this approach drastically curtails the memory and computational needs of LLMs. There is a clear movement from 16-bit down to more compact formats, including widely explored 4-bit schemes~\cite{gptq, awq}. Yet, despite the widespread adoption of post-training quantization in production LLMs, it generally yields sub-optimal performance.

Emerging research on 1-bit model designs, exemplified by BitNet~\cite{bitnet}, points to an encouraging path for cutting the operational costs of LLMs without degrading accuracy. Standard LLMs rely on 16-bit floating-point formats (FP16 or BF16), with matrix multiplication representing the core computational workload. Consequently, a substantial share of resource usage stems from floating-point multiplications and additions. Conversely, in BitNet, matrix operations are restricted to integer additions, which dramatically lowers the energy requirements for LLMs. Given that power efficiency constrains computational performance on many hardware platforms, such reductions in energy use can be leveraged to achieve increased computational speed.

Beyond compute resources, inference incurs significant expense due to the transfer of model weights from off-chip DRAM to faster, on-chip memories like SRAM. Attempts to expand SRAM in accelerators to improve throughput often result in much higher financial and energy costs than DRAM. 1-bit LLM architectures present a substantial advantage over full-precision counterparts by drastically decreasing the overall memory requirements—both in terms of storage capacity and memory bandwidth—which streamlines the process of moving weights and boosts inference speed and cost-effectiveness.

Here, we present a noteworthy 1-bit LLM extension, named \textbf{\text{BitNet b1.58}}, in which all parameters are ternary and take on values from \{-1, 0, 1\}. By incorporating 0 as an allowable parameter, in addition to the original values of 1-bit BitNet, the effective representation expands to 1.58 bits per parameter. \text{BitNet b1.58} maintains all computational benefits associated with the original BitNet, notably its matrix multiplication regime that nearly eliminates the need for multiplications, enabling efficient hardware realization. This variant exhibits the same favorable energy profile as its predecessor and outperforms traditional FP16 LLMs regarding memory efficiency, data throughput, and inference latency. Moreover, \text{BitNet b1.58} introduces two additional strengths. First, its expressive power is enhanced by enabling feature selection through the use of explicit zeros in parameter values, thereby improving the representational capacity over standard 1-bit models. Second, empirical results indicate that \text{BitNet b1.58}, with configurations starting from 3B parameters, is capable of equaling FP16 baselines in both perplexity and task-specific performance when controlling for training regime and architecture specifications.

\section{BitNet b1.58}

\text{BitNet b1.58} is constructed upon the BitNet framework, wherein the standard Transformer architecture substitutes \emph{nn.Linear} layers with \emph{BitLinear} modules. Training occurs from the beginning, utilizing weights quantized to 1.58 bits and activations quantized to 8 bits. This version incorporates several changes relative to the original BitNet, as detailed below.

\paragraph{Quantization Function.} The employed quantization procedure restricts the weights to values in $\{-1, 0, +1\}$ by utilizing an \emph{absmean} quantization function. Here, the weight matrix is normalized by its mean absolute value, and each entry is rounded to the closest integer within the allowable set:

\begin{equation}
    \widetilde{W} = \mathrm{RoundClip}(\frac{W}{\gamma + \epsilon}, -1, 1),
\end{equation}

\begin{equation}
    \mathrm{RoundClip}(x, a, b) = \max(a, \min(b, \mathrm{round}(x))),
\end{equation}

\begin{equation}
    \gamma = \frac{1}{nm}\sum_{ij} |W_{ij}|.
\end{equation}

For activations, the quantization technique aligns with that used in BitNet, with one key distinction: activations are not rescaled before the application of nonlinearities to the interval $[0, Q_b]$. Instead, scaling is performed to map all activation values for each token into the range $[-Q_b, Q_b]$, thereby eliminating zero-point quantization. This approach simplifies both software implementation and system-level optimization and was observed to have an insignificant impact on performance in our empirical studies.

\paragraph{LLaMA-alike Components.}

The LLaMA architecture~\cite{llama, llama2} has become the standard foundation for open-source LLM development. To align with the practices of the open-source community, our design of \text{BitNet b1.58} incorporates several features characteristic of LLaMA. Notably, RMSNorm~\cite{rmsnorm}, SwiGLU~\cite{swiglu}, and rotary embedding~\cite{rope} are included, and all biases are omitted. These choices facilitate the seamless incorporation of \text{BitNet b1.58} into widely used open-source frameworks such as Huggingface, vLLM~\cite{vllm}, and llama.cpp\footnote{\url{https://github.com/ggerganov/llama.cpp}}, requiring minimal adaptation effort.

\section{Results}

We conducted a comparative analysis between \text{BitNet b1.58} and our own FP16 \text{LLaMA LLM} reproduction across several model scales. Both model families were pre-trained on the RedPajama dataset~\cite{redpajama} using a budget of 100 billion tokens to maintain experimental consistency. For evaluation, we assessed their zero-shot capabilities on diverse language understanding benchmarks: ARC-Easy~\cite{arc}, ARC-Challenge~\cite{arc}, Hellaswag~\cite{hellaswag}, Winogrande~\cite{winoGrande}, PIQA~\cite{piqa}, OpenbookQA~\cite{openbookqa}, and BoolQ~\cite{boolq}. Additionally, validation perplexity was measured using WikiText2~\cite{wikitext2} and C4~\cite{c4} datasets.

To further compare \text{LLaMA LLM} and \text{BitNet b1.58}, we quantified runtime GPU memory consumption and inference latency. These metrics were obtained with the highly optimized FasterTransformer framework\footnote{\url{https://github.com/NVIDIA/FasterTransformer}}, which delivers efficient LLM inference on GPUs. For \text{BitNet b1.58}, the Ladder~\cite{ladder} 2-bit kernel was incorporated. We specifically tracked the duration per generated token, as this constitutes the primary inference expense.

Table~\ref{tab:ppl} presents a summary of perplexity values and resource costs for both \text{BitNet b1.58} and \text{LLaMA LLM}. The data indicate that, beginning at the 3B parameter mark, \text{BitNet b1.58} achieves comparable perplexity to the full-precision \text{LLaMA LLM}, yet with 2.71-fold improvement in speed and 3.55-fold reduction in GPU memory use. Notably, the 3.9B variant of \text{BitNet b1.58} exhibits a 2.4-fold speedup, 3.32-fold decrease in memory requirements, and surpasses the 3B \text{LLaMA LLM} in performance.

A detailed breakdown of zero-shot task accuracy is given in Table~\ref{tab:zero-shot}. Evaluations adhered to the procedures outlined in the \emph{lm-evaluation-harness}\footnote{\url{https://github.com/EleutherAI/lm-evaluation-harness}} toolkit. The results indicate that as model size increases, the accuracy gap between \text{BitNet b1.58} and \text{LLaMA LLM} diminishes. Importantly, parity with the full-precision baseline is reached by \text{BitNet b1.58} at 3B parameters. Echoing the findings for perplexity, downstream task performance demonstrates that \text{BitNet b1.58} 3.9B yields superior results relative to \text{LLaMA LLM} 3B, while maintaining lower latency and memory requirements. Collectively, these results establish \text{BitNet b1.58} as a Pareto improvement over contemporary LLM architectures.

\paragraph{Memory and Latency}

We extended our experiments to larger model configurations, including 7B, 13B, and 70B parameters, and systematically assessed the associated computational cost. Figure~\ref{fig:latency-memory-energy} presents the observed patterns in both latency and memory requirements, revealing a positive correlation between speed-up and model scaling. Notably, \text{BitNet b1.58} at 70B demonstrates a speed advantage, operating at 4.1 times the speed of the \text{LLaMA LLM} reference. This acceleration can be attributed to the fact that the execution time for \emph{nn.Linear} operations expands as the model parameters increase. Memory usage displays a parallel progression, given that embeddings retain full-precision format, which becomes a proportionally lesser component of total memory in larger architectures. The benchmarks for latency and memory are conducted with a 2-bit kernel; hence, additional efficiency gains may be possible with further optimizations.

\paragraph{Energy}

We further approximated the energy consumption tied to arithmetic operations for both \text{BitNet b1.58} and \text{LLaMA LLM}. Our evaluation centers primarily on matrix multiplication, as it remains the predominant contributor to large language model operational cost. Figure~\ref{fig:energy} details the breakdown of energy use, highlighting that most computations in \text{BitNet b1.58} involve INT8 additions, whereas \text{LLaMA LLM} operations are split between FP16 additions and multiplications. Based on the energy estimates described in~\cite{energycost, pokebnn}, \text{BitNet b1.58} reduces the energy consumed by arithmetic operations in matrix multiplication by a factor of 71.4 on 7nm hardware. In addition, we present end-to-end energy consumption results for processing 512 tokens per model. The data indicate that as model size increases, the energy efficiency gap between \text{BitNet b1.58} and the FP16-based \text{LLaMA LLM} widens, favoring the former. This improvement arises because the share of \emph{nn.Linear} computations increases with scale, while overhead from other subsystems diminishes in proportion for larger models.

\paragraph{Throughput}

To assess throughput, we benchmarked \text{BitNet b1.58} and the 70B-parameter \text{LLaMA LLM} on two 80GB A100 GPUs, employing pipeline parallelism~\cite{gpipe} to facilitate the execution of \text{LLaMA LLM} 70B across these devices. We gradually increased the batch size until reaching the memory constraints of the GPU, while maintaining a fixed sequence length of 512. As detailed in Table~\ref{tab:throughput}, \text{BitNet b1.58} 70B is capable of accommodating up to 11 times larger batches compared to \text{LLaMA LLM}, leading to an 8.9-fold increase in throughput.

\textbf{\text{BitNet b1.58} introduces a revised scaling law for balancing model capability against inference costs.} Referring to Figures~\ref{fig:latency-memory-energy} and \ref{fig:energy}, we observe an equivalence between model sizes in 1.58-bit and 16-bit quantizations, summarized as follows:

\begin{itemize}
    \item BitNet b1.58 at 13B surpasses a 3B FP16 LLM regarding latency, memory footprint, and energy demand.
    \item The 30B BitNet b1.58 demonstrates better efficiency in latency, memory, and energy usage than a 7B FP16 LLM.
    \item BitNet b1.58 at 70B achieves higher efficiency with respect to latency, memory use, and energy than the 13B FP16 LLM.
\end{itemize}

\paragraph{Training with 2T Tokens}

The total number of training tokens significantly influences LLM development. To evaluate how \text{BitNet b1.58} scales with token counts, we conducted pretraining of a \text{BitNet b1.58} model on 2T tokens, using the StableLM-3B~\cite{StableLM-3B-4E1T} data pipeline as guidance, representative of a leading open-source 3B model. We assessed both models on a composite benchmark including Winogrande~\cite{winoGrande}, PIQA~\cite{piqa}, SciQ~\cite{sciq}, LAMBADA~\cite{lambada}, and ARC-easy~\cite{arc}. Zero-shot performance metrics are summarized in Table~\ref{tab:2t}; for benchmarks reporting both accuracy and normalized accuracy, we averaged these values. The StableLM 3b results for 2T tokens were sourced directly from its technical documentation. Our results reveal that \text{BitNet b1.58} consistently achieves higher performance across all tested tasks, demonstrating strong generalization in 1.58-bit LLMs.

\section{Discussion and Future Work}

\textbf{1-bit Mixture-of-Experts (MoE) LLMs}

The Mixture-of-Experts (MoE) framework has emerged as a resource-efficient strategy for deploying LLMs, primarily due to its ability to decrease computational FLOPs. Nevertheless, MoE architectures often incur considerable memory demands and significant inter-chip communication costs, which present challenges for widespread adoption. Employing 1.58-bit LLMs offers viable solutions to these issues. Reducing the memory requirements enables models to operate with fewer hardware devices. Additionally, it decreases the communication overhead associated with activation transfer between networked components. Ideally, should model size permit, situating the entire MoE network on a single chip would eliminate such overhead altogether.

\textbf{Native Support of Long Sequence in LLMs}

Supporting extended sequence lengths is now a foundational requirement for contemporary LLMs. A central obstacle for inference on long sequences lies in the substantial memory usage driven by KV caches. The advent of \text{BitNet b1.58} marks meaningful progress toward inherent compatibility with long-sequence processing, as activation size is reduced from 16 bits to 8 bits, enabling twice the context window within identical hardware constraints. Furthermore, for 1.58-bit LLMs, these activations could be compressed—without information loss—down to 4 bits or even less, a direction we propose for subsequent investigation.

\textbf{LLMs on Edge and Mobile}

Deploying 1.58-bit LLMs presents substantial advantages for language modeling tasks on resource-constrained platforms such as edge and mobile devices. Due to stringent memory and processing limitations, the scope and efficiency of LLMs on these devices have been traditionally constrained. Lowered memory consumption and energy requirements with 1.58-bit LLMs now allow effective deployment on such platforms, unlocking a range of novel applications and expanding their functionality. Importantly, these quantized LLMs exhibit strong compatibility with CPU architectures prevalent in edge and mobile systems, meaning \text{BitNet b1.58} can operate efficiently, thus broadening device capabilities and improving overall system performance.

\textbf{New Hardware for 1-bit LLMs}

Initiatives like~Groq\footnote{\url{https://groq.com/}} illustrate the substantial promise of purpose-built hardware, including LPUs, tailored for LLM workloads. We advocate further for the conceptualization and development of hardware and system-level architectures explicitly optimized for 1-bit LLMs, leveraging the computation paradigm shift introduced by BitNet~\cite{bitnet}.